\documentclass[11pt]{article}
\usepackage[margin=1in]{geometry}
\usepackage{amsmath,amssymb,amsthm}
\usepackage{hyperref}

\newtheorem{theorem}{Theorem}
\newtheorem{lemma}[theorem]{Lemma}

\title{Proof of All Three OEIS~A397048 Conjectures via Falling-Factorial Basis Shifting}
\author{Salem Levinstein}
\date{November 26, 2026}

\begin{document}
\maketitle

\begin{abstract}
We prove all three conjectures stated by Werner Schulte in \texttt{OEIS~A397048} for the triangle $T(n, k) = \frac{(-1)^k}{k!} \sum_{i=0}^{k} (-1)^i \binom{k}{i} (n + i)^n$, together with a general base-point $c \ge 0$ Stirling shift identity. Expressing $(x + c)^n$ in the falling-factorial basis $\{(x)_k\}$ and applying binomial orthogonality shows that $T_c(n, k) = \sum_{j=0}^{n} \binom{n}{j} c^{n - j} S(j, k) \in \mathbb{N}$ is the $k$-th falling-factorial coordinate of $(x + c)^n$. This immediately yields the row-sum identity $\sum_{k=0}^{n} T(n, k) = \mathrm{A134980}(n)$ (Conjecture~2) and the falling-factorial inversion $\sum_{i=0}^{k} T(n, i)\, k!/(k - i)! = (n + k)^n$ for every $k \ge 0$ (Conjecture~3). Shifting $x \mapsto x - c$ and multiplying by $(x)_c$ gives $x^n (x)_c = \sum_k T_c(n, k) (x)_{c + k}$; expanding the left-hand side via Stirling numbers of the first and second kind proves $T_c(n, k) = \sum_{m=0}^{c} s(c, m) S(n + m, c + k)$, which specializes at $c = n$ to Conjecture~1. All theorems are formally verified in Lean~4.
\end{abstract}

\section{Introduction}

In \texttt{OEIS~A397048}~\cite{Schulte2026}, Werner Schulte defines the triangle
\[
  T(n, k) = \frac{(-1)^k}{k!} \sum_{i=0}^{k} (-1)^i \binom{k}{i} (n + i)^n \qquad (0 \le k \le n)
\]
and records three conjectures, where $s(a, b) = \mathrm{A048994}(a, b)$ are the signed Stirling numbers of the first kind, $S(a, b) = \mathrm{A048993}(a, b)$ are the Stirling numbers of the second kind, and $B_j = \sum_{k=0}^{j} S(j, k) = \mathrm{A000110}(j)$ are the Bell numbers:
\begin{enumerate}
  \item $T(n, k) = \sum_{i=0}^{n - k} s(n, k + i)\, S(n + k + i, n + k)$;
  \item $\sum_{k=0}^{n} T(n, k) = \mathrm{A134980}(n) := \sum_{j=0}^{n} \binom{n}{j} n^{n - j} B_j$;
  \item $\sum_{i=0}^{k} T(n, i)\, \dfrac{k!}{(k - i)!} = (n + k)^n$.
\end{enumerate}
For any integer base point $c \ge 0$ and integers $n, k \ge 0$, set
\[
  T_c(n, k) = \frac{1}{k!} \sum_{i=0}^{k} (-1)^{k - i} \binom{k}{i} (c + i)^n,
\]
so that $T(n, k) = T_n(n, k)$. We prove all three conjectures and a general base-point-$c$ generalization of Conjecture~1.

\section{Falling-Factorial Preliminaries}

For $k \in \mathbb{N}$, define the falling-factorial polynomial $(x)_k \in \mathbb{Q}[x]$ by $(x)_0 = 1$ and $(x)_{k + 1} = (x)_k (x - k)$. For $m \in \mathbb{N}$, $(m)_k = m! / (m - k)! = k! \binom{m}{k}$ when $k \le m$ and $(m)_k = 0$ when $k > m$.

\begin{lemma}\label{lem:prelim}
Let $s(n, k) \in \mathbb{Z}$ and $S(n, k) \in \mathbb{N}$ be defined by $s(0, 0) = S(0, 0) = 1$, $s(0, k) = S(0, k) = 0$ for $k \ge 1$, $s(n + 1, 0) = S(n + 1, 0) = 0$, and $s(n + 1, k) = s(n, k - 1) - n\, s(n, k)$, $S(n + 1, k) = S(n, k - 1) + k\, S(n, k)$ for $k \ge 1$. Then:
\begin{enumerate}
  \item If $\sum_{k=0}^{N} a_k (x)_k = \sum_{k=0}^{N} b_k (x)_k$ in $\mathbb{Q}[x]$, then $a_k = b_k$ for all $0 \le k \le N$.
  \item $(x)_c (x - c)_k = (x)_{c + k}$ for all $c, k \in \mathbb{N}$.
  \item $(x)_n = \sum_{k=0}^{n} s(n, k) x^k$ and $x^n = \sum_{k=0}^{n} S(n, k) (x)_k$, with $s(n, k) = S(n, k) = 0$ for $k > n$.
  \item $\sum_{i=0}^{k} (-1)^{k - i} \binom{k}{i} (i)_m = k!\, \delta_{k, m}$ for all $k, m \in \mathbb{N}$.
\end{enumerate}
\end{lemma}

\begin{proof}
Part~(1) follows by strong induction on $k \le N$ upon evaluating at $x = k$, since $(k)_j = 0$ for $j > k$ and $(k)_k = k! \ne 0$. Parts~(2) and~(3) follow immediately by induction on $k$ and $n$ using $(x)_{n + 1} = (x)_n(x - n)$ and $x(x)_k = (x)_{k + 1} + k(x)_k$. For~(4), if $m > k$ then $(i)_m = 0$ for every $0 \le i \le k$; if $m \le k$, then $(i)_m = m!\binom{i}{m}$ vanishes for $i < m$, and re-indexing $i = m + r$ with $\binom{k}{m + r}\binom{m + r}{m} = \binom{k}{m}\binom{k - m}{r}$ gives $m!\binom{k}{m}\sum_{r=0}^{k-m}(-1)^{(k-m)-r}\binom{k-m}{r} = m!\binom{k}{m}(1 - 1)^{k - m} = k!\,\delta_{k, m}$.
\end{proof}

\section{Main Theorems}

\begin{theorem}\label{thm:exp}
For all integers $c, n, k \ge 0$,
\[
  T_c(n, k) = \sum_{j=0}^{n} \binom{n}{j} c^{n - j} S(j, k) \in \mathbb{N}.
\]
In particular, $T_c(n, k) = 0$ for $k > n$, and $(x + c)^n = \sum_{k=0}^{n} T_c(n, k)\, (x)_k$ in $\mathbb{Q}[x]$.
\end{theorem}

\begin{proof}
Set $U_c(n, k) = \sum_{j=0}^{n} \binom{n}{j} c^{n - j} S(j, k)$; then $U_c(n, k) = 0$ for $k > n$. By the binomial theorem and Lemma~\ref{lem:prelim}(3),
\[
  (x + c)^n = \sum_{j=0}^{n} \binom{n}{j} c^{n - j} x^j = \sum_{j=0}^{n} \binom{n}{j} c^{n - j} \sum_{m=0}^{n} S(j, m)\, (x)_m = \sum_{m=0}^{n} U_c(n, m)\, (x)_m.
\]
Fix $k \ge 0$ and let $N = \max(n, k)$. Evaluating at $x = i$ ($0 \le i \le k$), multiplying by $(-1)^{k - i} \binom{k}{i}$, summing over $i$, and applying Lemma~\ref{lem:prelim}(4) gives
\[
  k!\, T_c(n, k) = \sum_{m=0}^{N} U_c(n, m) \sum_{i=0}^{k} (-1)^{k - i} \binom{k}{i} (i)_m = k!\, U_c(n, k).
\]
Dividing by $k! \ne 0$ yields $T_c(n, k) = U_c(n, k)$ and the claimed expansion.
\end{proof}

\begin{theorem}\label{thm:conjs}
For all integers $c, n, k \ge 0$:
\begin{enumerate}
  \item \emph{(General shift theorem and Conjecture~1.)} If $k \le n$, then
  \[
    T_c(n, k) = \sum_{m=0}^{c} s(c, m)\, S(n + m, c + k).
  \]
  In particular, for $c = n$ and $0 \le k \le n$,
  \[
    T(n, k) = \sum_{i=0}^{n - k} s(n, k + i)\, S(n + k + i, n + k).
  \]
  \item \emph{(Conjecture~2.)} $\sum_{k=0}^{n} T(n, k) = \sum_{j=0}^{n} \binom{n}{j} n^{n - j} B_j = \mathrm{A134980}(n)$.
  \item \emph{(Conjecture~3.)} For every $n, k \ge 0$, $\sum_{i=0}^{k} T(n, i)\, \dfrac{k!}{(k - i)!} = (n + k)^n$.
\end{enumerate}
\end{theorem}

\begin{proof}
For~(1), substitute $x \mapsto x - c$ in Theorem~\ref{thm:exp} to obtain $x^n = \sum_{j=0}^{n} T_c(n, j)\, (x - c)_j$, multiply by $(x)_c$, and apply Lemma~\ref{lem:prelim}(2):
\[
  x^n (x)_c = \sum_{j=0}^{n} T_c(n, j)\, (x)_{c + j}.
\]
On the other hand, Lemma~\ref{lem:prelim}(3) yields
\[
  x^n (x)_c = \sum_{m=0}^{c} s(c, m)\, x^{n + m} = \sum_{K=0}^{c + n} \left(\sum_{m=0}^{c} s(c, m)\, S(n + m, K)\right) (x)_K.
\]
Comparing the coefficients of $(x)_{c + k}$ ($0 \le k \le n$) via Lemma~\ref{lem:prelim}(1) proves
\[
  T_c(n, k) = \sum_{m=0}^{c} s(c, m)\, S(n + m, c + k).
\]
At $c = n$, terms with $m < k$ vanish since $S(n + m, n + k) = 0$, and setting $m = k + i$ ($0 \le i \le n - k$) yields Conjecture~1.

For~(2), summing Theorem~\ref{thm:exp} at $c = n$ over $0 \le k \le n$ and exchanging finite sums (using $S(j, k) = 0$ for $k > j$) gives
\[
  \sum_{k=0}^{n} T(n, k) = \sum_{j=0}^{n} \binom{n}{j} n^{n - j} \sum_{k=0}^{j} S(j, k) = \mathrm{A134980}(n).
\]

For~(3), set $N = \max(n, k)$ and evaluate $(x + n)^n = \sum_{i=0}^{N} T(n, i)\, (x)_i$ at $x = k$: terms with $i > k$ vanish, and $(k)_i = k! / (k - i)!$ for $0 \le i \le k$.
\end{proof}

\section{Formal Verification}

Lemma~\ref{lem:prelim}, Theorem~\ref{thm:exp}, and Theorem~\ref{thm:conjs} have been formally verified in Lean~4 with Mathlib (module \texttt{Conjectures.SchulteA397048}) using only the standard axioms (no \texttt{sorry}, custom axioms, or \texttt{native\_decide}). An independent Python verifier (\texttt{verify.py}) confirms $4{,}424$ instances across Conjectures~1--3 ($n \le 40$) and the general shift theorem ($c, n \le 14$).

\begin{thebibliography}{9}
\bibitem{Schulte2026}
W.~Schulte,
\emph{OEIS Sequence A397048},
The On-Line Encyclopedia of Integer Sequences, June~14, 2026 (rev.~34),
\url{https://oeis.org/A397048}.
\end{thebibliography}

\end{document}
